% Ad Atticum 2.16 --- headnote
% ad-atticum-02-16, end of April or early May 59 BC, written at Formiae

\textit{Cicero to Atticus, written from the Formian villa
at the very end of April or beginning of May 59 BC. Caesar's
agrarian law has been carried at the start of the year; the
news that has reached Cicero on his couch is that the
Campanian land --- previously kept as a state asset and
exempted from the bill --- is to be added to the distribution
under a second law. Cicero turns it over the night through
and writes \textsection 1 the next day in self-consolation:
the Campanian land cannot support more than five thousand
men, the rest of the urban populace will be alienated from
the triumvirs by it, and the abolition of the tolls of Italy
plus this distribution will leave only the inheritance-tax
of one twentieth as Roman state revenue --- which will be
``bound to perish in one little public meeting.''}

\textit{\textsection 2 is the most exuberant single
paragraph of the letter corpus on Pompey's transformation
from constitutionalist to agent of the triumvirate. The
Greek tag from Sophocles or some lost tragedian (``he blows
now no longer with little pipes but with savage bellows
without the muzzle'') sets the key. Cicero impersonates
Pompey through the standard sophistic defences he has been
giving (``he approved Caesar's laws,'' ``the agrarian law
pleased him,'' ``about Bibulus watching the sky\ldots''),
then turns and addresses him as ``Sampsiceramus'' --- the
nickname Cicero adopts for Pompey throughout 59, after the
petty Syrian dynast Pompey had reduced --- demanding to know
how he will defend taking the Campanian land too. The
imagined Pompey replies: ``I shall hold you down, crushed,
by Caesar's army.'' Cicero answers that what holds him down
is not the army but the ingratitude of the so-called good
men, who have given him no thanks even in conversation.}

\textit{The letter resolves with the famous philosophical
choice (\textsection 3): Dicaearchus (Atticus's man, the
Peripatetic of the active life) and Theophrastus (Cicero's
man, of the contemplative life) are at odds. Cicero has,
he says, given Dicaearchus full measure; now he turns to
Theophrastus and resolves to ``bend, our dear Titus, to
those splendid pursuits, and go back at last to the place
from which we ought never to have departed.'' The famous
turn from politics to philosophy at the doorstep of exile
is here in skeleton.}

\textit{The closing paragraph (\textsection 4) is on
Quintus, whose latest letter is a Chimaera ``lion in front,
snake behind'' --- lamenting his Asian stay yet asking
Cicero to correct and publish his annals --- and on a
practical question of the Asian transport-toll which the
publicans (Pompey's allies) are pressing.}
